\documentclass[12pt]{article}

\usepackage{amsmath, amssymb}
\usepackage{booktabs}
\usepackage{geometry}
\usepackage{setspace}

\geometry{margin=1in}
\setstretch{1.2}

\title{Lecture 18 Scribe\\
Marginal PMF, Correlation, Covariance, Joint Gaussian Random Variables}
\author{}
\date{}

\begin{document}

\maketitle

\section*{1. Marginal PMF}

\subsection*{1.1 Definition}

For discrete random variables $X$ and $Y$, the joint PMF is defined as

\[
p_{XY}(x,y) = P(X = x, Y = y)
\]

This function gives the probability that both variables simultaneously take particular values.

The marginal PMF of each variable is obtained by summing over the other variable:

\[
p_X(x) = \sum_y p_{XY}(x,y)
\]

\[
p_Y(y) = \sum_x p_{XY}(x,y)
\]

\subsubsection*{Reasoning Behind the Definition}

The joint PMF describes the combined behavior of two variables.

However, sometimes we want to know the distribution of only one variable regardless of the value of the other variable.

To obtain this, we sum all probabilities involving that value of the variable across every possible value of the other variable.

\textbf{Key idea:} The marginal PMF of one variable is obtained by summing the joint PMF over all possible values of the other variable.

\section*{1.2 Marginal PMF from a Joint PMF Table}

When the joint PMF is given in a table, marginal PMFs can be computed by summing rows and columns.

\begin{center}
\begin{tabular}{cccc}
\toprule
$p_{XY}(x,y)$ & $y_1$ & $y_2$ & $y_3$ \\
\midrule
$x_1$ & $p_{11}$ & $p_{12}$ & $p_{13}$ \\
$x_2$ & $p_{21}$ & $p_{22}$ & $p_{23}$ \\
$x_3$ & $p_{31}$ & $p_{32}$ & $p_{33}$ \\
\bottomrule
\end{tabular}
\end{center}

From this structure

\[
p_X(x_i) = \sum_j p_{ij}
\]

\[
p_Y(y_j) = \sum_i p_{ij}
\]

\textbf{Reasoning}

\begin{itemize}
\item Row sums $\rightarrow$ marginal distribution of $X$
\item Column sums $\rightarrow$ marginal distribution of $Y$
\end{itemize}

Also the marginal PMFs must satisfy

\[
\sum_i p_X(x_i)=1
\]

\[
\sum_j p_Y(y_j)=1
\]

so they remain valid probability distributions.

\section*{1.3 Example: Computing Marginal PMFs}

Given joint PMF

\begin{center}
\begin{tabular}{ccc}
\toprule
$p_{XY}(x,y)$ & $y=0$ & $y=1$ \\
\midrule
$x=0$ & $\frac{1}{8}$ & $\frac{1}{4}$ \\
$x=1$ & $\frac{3}{8}$ & $\frac{1}{4}$ \\
\bottomrule
\end{tabular}
\end{center}

\subsection*{Step 1: Compute $p_X(x)$}

\[
p_X(0)=\frac{1}{8}+\frac{1}{4}
\]

\[
=\frac{1}{8}+\frac{2}{8}=\frac{3}{8}
\]

\[
p_X(1)=\frac{3}{8}+\frac{1}{4}
\]

\[
=\frac{3}{8}+\frac{2}{8}=\frac{5}{8}
\]

Thus

\[
p_X(x)=
\begin{cases}
3/8 & x=0 \\
5/8 & x=1
\end{cases}
\]

\subsection*{Step 2: Compute $p_Y(y)$}

\[
p_Y(0)=\frac{1}{8}+\frac{3}{8}=\frac{4}{8}=\frac{1}{2}
\]

\[
p_Y(1)=\frac{1}{4}+\frac{1}{4}=\frac{1}{2}
\]

Thus

\[
p_Y(y)=
\begin{cases}
1/2 & y=0 \\
1/2 & y=1
\end{cases}
\]

\section*{2. Correlation}

A basic way to measure the joint behavior of two random variables is through the product moment

\[
R_{XY}=E[XY]
\]

\textbf{Reasoning}

This quantity represents the average value of the product $XY$.

However:

\begin{itemize}
\item $R_{XY}$ depends on the scale of $X$ and $Y$
\item It is computed about the origin rather than the means
\end{itemize}

Thus it mixes together mean values, spread, and dependence.

If

\[
R_{XY}=0
\]

then the variables are orthogonal about the origin.

To obtain a better measure of dependence we center the variables. This leads to covariance.

\section*{3. Covariance}

\subsection*{3.1 Definition}

\[
Cov(X,Y)=E[(X-\mu_X)(Y-\mu_Y)]
\]

where

\[
\mu_X=E[X], \quad \mu_Y=E[Y]
\]

\textbf{Interpretation}

\begin{itemize}
\item $Cov(X,Y)>0$ : positive covariance
\item $Cov(X,Y)<0$ : negative covariance
\item $Cov(X,Y)=0$ : no linear co-variation
\end{itemize}

\subsection*{3.2 Algebraic Form}

\[
Cov(X,Y)=E[(X-\mu_X)(Y-\mu_Y)]
\]

Expand

\[
=E[XY-X\mu_Y-Y\mu_X+\mu_X\mu_Y]
\]

Using linearity of expectation

\[
=E[XY]-\mu_YE[X]-\mu_XE[Y]+\mu_X\mu_Y
\]

Since

\[
\mu_X=E[X], \quad \mu_Y=E[Y]
\]

we obtain

\[
Cov(X,Y)=E[XY]-E[X]E[Y]
\]

\section*{3.3 Properties of Covariance}

\textbf{1. Symmetry}

\[
Cov(X,Y)=Cov(Y,X)
\]

\textbf{2. Variance as Special Case}

\[
Cov(X,X)=Var(X)
\]

\textbf{3. Linearity}

\[
Cov(aX+b,cY+d)=ac \, Cov(X,Y)
\]

\textbf{4. Independence}

If $X$ and $Y$ are independent

\[
Cov(X,Y)=0
\]

However zero covariance does not necessarily imply independence.

\section*{4. Example}

Given

\[
Y=-2X+3
\]

Find $Cov(X,Y)$.

\subsection*{Step 1}

\[
Cov(X,Y)=E[XY]-E[X]E[Y]
\]

\subsection*{Step 2}

\[
XY=X(-2X+3)
\]

\[
=-2X^2+3X
\]

\subsection*{Step 3}

\[
E[XY]=-2E[X^2]+3E[X]
\]

\[
E[Y]=-2E[X]+3
\]

After simplification

\[
Cov(X,Y)=-2Var(X)
\]

This indicates a negative relationship between $X$ and $Y$.

\section*{5. Pearson Correlation Coefficient}

To remove units we normalize covariance.

\[
\rho_{XY}=\frac{Cov(X,Y)}{\sqrt{Var(X)Var(Y)}}
\]

or

\[
\rho_{XY}=\frac{Cov(X,Y)}{\sigma_X\sigma_Y}
\]

Range

\[
-1\le \rho_{XY}\le 1
\]

Interpretation

\begin{itemize}
\item $\rho_{XY}>0$ : positive correlation
\item $\rho_{XY}<0$ : negative correlation
\item $\rho_{XY}=0$ : no linear correlation
\end{itemize}

\section*{6. Jointly Gaussian Random Variables}

A pair $(X,Y)$ is jointly Gaussian if

\begin{itemize}
\item the joint density has a bivariate Gaussian form
\item the marginal distributions are Gaussian
\end{itemize}

\subsection*{Joint Gaussian PDF}

\[
f_{XY}(x,y)=
\frac{1}{2\pi\sigma_X\sigma_Y\sqrt{1-\rho^2}}
\exp
\left(
-\frac{1}{2(1-\rho^2)}
\left[
\left(\frac{x-\mu_X}{\sigma_X}\right)^2
+
\left(\frac{y-\mu_Y}{\sigma_Y}\right)^2
-
2\rho
\left(\frac{x-\mu_X}{\sigma_X}\right)
\left(\frac{y-\mu_Y}{\sigma_Y}\right)
\right]
\right)
\]

where

\begin{itemize}
\item $\mu_X,\mu_Y$ are the means
\item $\sigma_X,\sigma_Y$ are standard deviations
\item $\rho$ is the correlation coefficient
\end{itemize}

\subsection*{Interpretation Example}

Example variables

\begin{itemize}
\item $X$: Atmospheric CO$_2$ concentration (ppm)
\item $Y$: Global surface temperature anomaly ($^\circ$C)
\end{itemize}

If

\[
\rho>0
\]

then higher CO$_2$ levels tend to occur with higher temperatures.

\end{document}